% Finite-bit decoder theorem and proof, 8 October 2026.
\subsection{Maintaining fixed-rank attention through exact decoding}
\label{sec:momentdecoder}

This subsection proves the decoder of Theorem~\ref{thm:main-decoder}.
An index that accepts already computed keys leaves their construction
outside its cost.  For a fixed key-map rank, polynomial moments also
permit the construction of the new keys and values.  We first give a
deterministic finite-precision evaluator.  We then use certified
probability intervals to sample exactly, with rare precision increases
paid for by their probability.  The maintained Taylor feature
sums~\cite{katharopoulos2020linear,heinsen2026taylor} and exact
sampling by interval refinement~\cite{han1997interval,brassard2019remote}
are established; Section~\ref{sec:prior-main} says what we take from
them.  The claims here concern finite arithmetic, a uniform depth error
bound, maintenance while precision increases, and the exact original
autoregressive law.

\paragraph{Model.}
The vocabulary has size $V$, the residual width is $n$, and there are
$D$ causal attention blocks.  Each block first applies
\[
 N_\varepsilon(h)=
 \frac{h}{\sqrt{\varepsilon+\|h\|_2^2/n}},\qquad \varepsilon>0,
\]
then affine query, key, and value maps.  Its attention scores are
$\beta\langle q_t,k_j\rangle/n$, $j\le t$.  The attention update
is added to the residual stream.  A second normalization followed by
an affine map and coordinatewise tanh gives the second additive
update.  Either additive update may be multiplied by a fixed dyadic
coefficient $\alpha\in[0,1]$.  Convex residual mixtures with
coefficients in $[0,1]$ are also allowed.
An additional fixed affine attention output map can be included.
The output consists of $V$ affine logits followed by softmax; the
logits are promised to lie in $[-1,1]$.  Input tokens use a fixed
dyadic embedding table with entries of magnitude at most $s$.

All model parameters are dyadic with encoding length at most $b$,
using binary integer numerators and denominators; in particular,
their denominators divide $2^b$.
The row sums of absolute values and the bias bounds are at most
$s\ge1$.  Every positive normalization parameter is at least
$2^{-b}$.  The temperature $\beta\ge0$ is dyadic of at most $b$
bits, or $\beta=\sqrt n$.  Each layer uses one fixed affine key
map at every position, and its linear part has rank at most the
fixed integer $r$.  Thus there is one common known affine key
space per layer containing the keys at all positions.  Separate
rank bounds for position-dependent maps, such as positional
rotations that enlarge the joint key span, do not imply this
joint-space condition.  The affine key space is computed from the
fixed weights, independently of the token history.

The statement also applies to a bounded number of heads, with a
separate rank promise for each head.  More generally, the number of
heads can be included as another polynomial parameter.  We state
the formulas for one head to keep notation short.

\begin{theorem}[Exact decoding with a fixed-rank key map]
\label{thm:momentdecoder}
For each fixed $r$, the model above has a uniform exact autoregressive
sampler.  Let $T_0\ge1$ be the prompt length and $T\ge T_0$ the
current context length.  There is a constructive polynomial $F_r$
such that prefill costs at most
\[
 T_0 F_r(n,D,V,b,s,\beta,\log(T_0+2)),
\]
and every subsequent token has expected bit cost at most
\[
 F_r(n,D,V,b,s,\beta,\log(T+2)).
\]
The expectation is conditional on every previous token history.
It includes the construction and insertion of the new keys and
values, random bits, and all precision refinement.  The sampler
uses no exact-real oracle.  At standard scaling $\beta=\sqrt n$,
if $s,D,V,b$ are polynomially bounded in $n$, the expected work is
$\operatorname{poly}_r(n,\log T)$.  The context dependence outside
the stored token history is polylogarithmic.

The same argument applies to a specified family of other
feedforward updates when fixed polynomial bounds for its range,
certified evaluation cost, and logarithm of its Lipschitz bound
are supplied as part of the theorem's assumptions.  The resulting
polynomial also depends on those fixed bounds.  The degree of
$F_r$ depends on $r$; the uniform polynomial dependence on width
requires a rank bound independent of width.
\end{theorem}

We prove the theorem after three lemmas: a finite moment index
(Lemma~\ref{lem:finitemomentindex}), uniform precision for a causal
replay (Lemma~\ref{lem:prefixstability}), and exact sampling from
certified replays (Lemma~\ref{lem:replayexactsampling}).

\begin{lemma}[A finite moment index]
\label{lem:finitemomentindex}
Fix $r$.  Let $z_j\in[-1,1]^r$ and $u_j\in[-1,1]^a$ be dyadic
records on a common grid $2^{-P}$, and let $t\in\mathbb R^r$
satisfy $\|t\|_1\le H$, where $H\ge1$ is an integer.  The
coordinates of $t$ are either rational with a common denominator,
or have the form $c_i/\sqrt n$ with rational $c_i$ having a common
denominator.  Write
\[
 A(t)=\frac{\sum_{j=1}^T e^{\langle t,z_j\rangle}u_j}
              {\sum_{j=1}^T e^{\langle t,z_j\rangle}}.
\]
For each $p\ge1$, an append-only deterministic index returns
certified intervals of width at most $2^{-p}$ for all coordinates
of $A(t)$.  Its update and query costs are polynomial, at fixed
$r$, in $a,H,P,p$, the rational coefficient lengths, $\log(T+2)$,
and, in the second case, $\log(n+2)$.
Its arithmetic has no error that accumulates with the number of
insertions.
\end{lemma}

\paragraph{Proof idea.}
Expand the exponential in the fixed number of key coordinates.
The index stores exact sums of monomials and value-weighted
monomials.  Their numerator lengths grow by $\log T$.  A uniform
Taylor error controls the normalized attention mean because its
average denominator is at least $e^{-H}$.

\begin{proof}
Put
\[
 q=p+2H+6,\qquad m=\lceil8(H+q+1)\rceil,\qquad
 \mathcal I_m=\{\nu\in\mathbb N^r:|\nu|\le m\}.
\]
Let $N_m=|\mathcal I_m|=\binom{m+r}{r}$.  The Taylor polynomial
$E_m(x)=\sum_{d=0}^m x^d/d!$ satisfies
\begin{equation}
 \sup_{|x|\le H}|E_m(x)-e^x|\le2^{-q}.
 \label{eq:momenttaylorerror}
\end{equation}
Indeed, the Lagrange remainder is at most
$e^H H^{m+1}/(m+1)!$.  The elementary bound
$N!\ge(N/e)^N$ follows by integrating $\log x$ below its sum.
Since $m+1\ge8H$ and $e<4$, the remainder is at most
$e^H2^{-(m+1)}\le2^{2H-m-1}\le2^{-q}$.

Store the exact rational sums
\[
 S_\nu=\sum_{j=1}^T z_j^\nu,\qquad
 M_{\nu,i}=\sum_{j=1}^T z_j^\nu u_{j,i},
 \qquad \nu\in\mathcal I_m,\quad i\le a.
\]
They can be stored as integer numerators over the one denominator
$2^{P(m+1)}$.  An insertion computes its monomials exactly and
adds the corresponding integer increments.  The absolute value
of every stored sum is at most $T$, so its numerator has
$O(P(m+1)+\log(T+2))$ bits.  No rounding occurs during these
additions.  Different dyadic coordinate grids can be put over one
common power-of-two denominator before forming the moments.

The multinomial theorem gives
\[
 E_m(\langle t,z\rangle)=
 \sum_{\nu\in\mathcal I_m}\frac{t^\nu z^\nu}{\nu!}.
\]
Consequently the truncated numerator and denominator are computed
from the stored moments by $O(aN_m)$ exact rational operations.
If the common denominator of the rational coordinates of $t$ is
$d_t$, all Taylor coefficients have a common denominator dividing
$d_t^m m!$.  This follows because $\nu!$ divides $|\nu|!$, which
divides $m!$.  Integer numerator lengths are polynomial in
$m,P$, the coefficient encoding lengths, and $\log T$.

Divide the true and truncated sums by $T$.  Their denominators
are $Z$ and $\widetilde Z$, and their coordinate numerators are
$N_i$ and $\widetilde N_i$.  Then
\[
 Z\ge e^{-H},\quad |N_i|\le Z,\quad
 |Z-\widetilde Z|\le\delta,\quad
 |N_i-\widetilde N_i|\le\delta,
 \qquad \delta=2^{-q}.
\]
The choice of $q$ makes $\delta\le Z/2$.  Thus
\[
 \left|\frac{\widetilde N_i}{\widetilde Z}
            -\frac{N_i}{Z}\right|
 \le\frac{2\delta}{Z-\delta}
 \le4\delta e^H
 \le2^{-p-4}.
\]
In particular the truncated denominator is positive.  A rational
division enclosure of width $2^{-p-2}$, enlarged by the displayed
Taylor error on both sides, has width less than $2^{-p}$.

At standard scaling the coordinates are $t_i=c_i/\sqrt n$.
Separate multiindices by the parity of $|\nu|$.  Every truncated
sum is then of the form $A+B/\sqrt n$ with rational $A,B$.
A common denominator is
$d_c^m n^{\lfloor m/2\rfloor}m!$, where $d_c$ is a common
denominator for the $c_i$.  Its encoding length is
$O(m\log n+m\log(m+2))$ plus the rational input lengths.
The ratio of two such expressions is evaluated by certified
bisection for $\sqrt n$.  The positive lower bound
$\widetilde Z\ge e^{-H}/2$ bounds the required guard precision.
If the rational coefficients have at most $L$ bits, precision
$O(L+H+p+\log(T+2))$ in the square root suffices.  This remains
polynomial.  If $n$ is a square, ordinary rational arithmetic
suffices.  Hence the temperature is exactly $\sqrt n$ throughout;
it has not been replaced by a floating-point constant.
\end{proof}

\paragraph{An explicit polynomial bound.}
The preceding construction can deliberately use schoolbook integer
arithmetic and recompute monomials.  With
\[
 L=C_r\{m[P+b+\log(n+m+D+V+2)]+\log(T+2)+1\},
\]
one may bound an entire token's deterministic finite-precision
evaluation by
\begin{equation}
 C_r(D+1)(n+V+1)^2(m+1)^{2r+4}L^4.
 \label{eq:momentexplicitcost}
\end{equation}
Here $b$ also bounds the fixed rational chart lengths up to a
factor depending on $r$, and $P$ includes the guard precision
described below.  To check \eqref{eq:momentexplicitcost}, note
that every integer formed for one token has $O(L)$ bits, so a
schoolbook product costs $O(L^2)$.  The layers and the output head
use $O((D+1)(n+V+1)^2)$ products in affine maps and $O((D+1)(n+V))$
certified scalar evaluations, each of cost $O(L^4)$ by
Lemma~\ref{lem:newnegativeexp} and rational bisection.  The moment
updates and queries use $O((D+1)(n+m+1)N_m)$ products, where
$N_m=\binom{m+r}{r}\le(m+1)^r$: each monomial and each Taylor
coefficient takes $O(m)$ products and meets $n$ value coordinates.  The total is
$O_r((D+1)(n+V+1)^2(m+1)^{r+1}L^4)$.  The displayed bound is
intentionally loose; it counts bit operations.

\begin{lemma}[Uniform stability over a causal prefix]
\label{lem:prefixstability}
For the model in Theorem~\ref{thm:momentdecoder}, a token history
of any length has a deterministic certified replay algorithm.  To
return all output probabilities and their cumulative sums within
absolute interval width $2^{-p}$, it uses working precision
\[
 P=p+O\!\left((D+1)[b+\log(n+s+\beta+2)]
                 +\log(D+V+2)\right).
\]
Its work per replayed token is bounded by a fixed polynomial in
$n,D,V,b,s,\beta,p,\log(T+2)$ at fixed $r$.  The stability guard
does not grow linearly with the context length.
\end{lemma}

\paragraph{Proof idea.}
An attention row is stable in the largest coordinate error, with
a bound independent of the number of keys.  Each layer reads only
the preceding layer's prefix.  Errors therefore propagate through
network depth.  Exact moment counters introduce no additional
recurrence over token positions.

\begin{proof}
For scores $a$ and values $v$, write $p(a)$ for softmax and
$F(a,v)=\sum_j p_j(a)v_j$.  Along a differentiable score path,
\[
 \frac{d}{dt}p_j=p_j\left(\dot a_j-\sum_i p_i\dot a_i\right).
\]
Summing absolute values shows that softmax is Lipschitz from
$\ell_\infty$ scores to $\ell_1$ probabilities with constant
two, for every number of positions.  If the value coordinates
are bounded by $K$, this gives
\begin{equation}
 \|F(a,v)-F(a',v')\|_\infty
 \le\|v-v'\|_\infty+2K\|a-a'\|_\infty.
 \label{eq:attentionstabilitymoment}
\end{equation}
For $|q_i|,|k_{j,i}|\le K$, perturbations bounded by
$\delta_q,\delta_k$ change any score by at most
$\beta K(\delta_q+\delta_k)$.  The same estimate follows by
interpolating between two bounded endpoint arrays, so it has no
extra product-error term.

Normalization has a global bound independent of the residual
radius.  Put $v=\varepsilon+\|h\|_2^2/n$.  Its Jacobian is
\[
 \frac{\partial N_i}{\partial h_j}
 =\frac{\mathbf1_{i=j}}{\sqrt v}
       -\frac{h_i h_j}{n v^{3/2}}.
\]
Using $|h_i|\le\sqrt{n(v-\varepsilon)}$ and
$\|h\|_1\le n\sqrt{v-\varepsilon}$, every absolute row sum
is at most $(1+\sqrt n)/\sqrt\varepsilon$.  Hence the
$\ell_\infty$ Lipschitz constant is at most
$(1+\sqrt n)2^{b/2}$.  Also $|N_i(h)|\le\sqrt n$.
Choose a known power of two
\[
 K\ge s(\lceil\sqrt n\rceil+1),\qquad
 K<2s(\lceil\sqrt n\rceil+1).
\]
Queries, keys, and values obtained from normalized inputs are
bounded by $K$.  Affine maps have Lipschitz bound $s$, tanh has
bound one, and residual addition has bound two on a product of
streams.  A fixed affine output map only introduces another
factor polynomial in $s$.  Thus the network is a composition of
$J=O(D+1)$ maps on prefix arrays, each with a common Lipschitz
bound $L_0\ge2$ satisfying
\[
 \log_2 L_0=O(b+\log(n+s+\beta+2)).
\]
These maps act on all positions with the maximum norm.  The bound
is independent of the number of positions by
\eqref{eq:attentionstabilitymoment}.

\paragraph{Induction over depth.}
If each primitive's computed output has local error at most
$\delta$, induction bounds the final error by
$J L_0^J\delta$.  Choose
\[
 \delta=2^{-P},\qquad
 P\ge p+\lceil\log_2(J+1)\rceil
       +J\lceil\log_2L_0\rceil+10.
\]
Known interval error bounds can be propagated with these same
constants, so this argument returns certified enclosures.  The
output head and its
cumulative probabilities are evaluated with the same margin.
Exponentials on bounded logit intervals and a positive softmax
denominator have elementary polynomial certified evaluation cost.
The $V$ sums and output coordinates are charged explicitly.

\paragraph{Rank-preserving rounding.}
A deterministic rounding operation first computes a certified
interval of width at most $2^{-P-3}$, then rounds its rational midpoint to the grid
$2^{-P}$.  Its local error is less than $2^{-P}$ after the fixed
guard margin is included.  It does not decide the exact nearest
grid point of an unknown real value.  Apply this rule to the output
of normalization, clipping the rounded value to the
rational interval $[-\lceil\sqrt n\rceil,\lceil\sqrt n\rceil]$,
and then compute
\[
 \widetilde k=W_K\widetilde N+b_K
\]
\emph{exactly}.  The result is dyadic on the grid $2^{-(P+b)}$
and lies exactly in $b_K+\operatorname{image}(W_K)$.  Independent
rounding of the coordinates of $\widetilde k$ is not used.  Query
and value affine products can be computed exactly in the same
way.  The completed residual updates are rounded to $2^{-P}$.
All these local rounding errors are included in $\delta$ by
requesting a fixed number of guard bits inside each primitive.

\paragraph{The moment chart.}
Use a bounded rational coordinate chart for the key image,
$k=c_0+A k_I$, with $|A_{ij}|\le2$.  The chart is computed once
from the weights by Lemma~\ref{lem:dynamicchart}, applied to a
column basis of $W_K$.  Put
$z_j=(k_j)_I/K$.  Then $z_j\in[-1,1]^r$, and the common score
offset $\beta\langle q,c_0\rangle/n$ cancels.  The remaining
coefficient vector is
\[
 t=\frac{\beta K}{n}q^T A,\qquad
 \|t\|_1\le2r\beta K^2.
\]
Take $H=\lceil\max\{1,2r\beta K^2\}\rceil$, using a rational
upper bound for $\sqrt n$ if necessary.  Scale values by $K$ and
apply Lemma~\ref{lem:finitemomentindex}; restoring that scale
needs only $O(\log(K+2))$ additional requested precision bits.
The key and value grids are common throughout a precision phase,
so the moment additions remain exact.  The common determinant
denominator in $A$ gives rational coefficient lengths
$O_r(P+b+\log(n+2))$; an $n$-term projection adds $O(\log n)$
numerator bits.

\paragraph{Causal order.}
Process a prefix in increasing position order.  At a layer and
position $t$, first form its query, key and value from the current
preceding-layer residual.  Append the current key and value to
that layer's moment index \emph{before} querying it, so the
attention range is exactly $j\le t$.  Complete the residual and
feedforward operations, then move to the next layer.  This is
the same deterministic computation whether performed during
prefill, an online append, or a replay.

To see why no token-count error factor is hidden, regard each
layer's exact and approximate states as arrays over all positions
in the current prefix.  Every attention input belongs to the
preceding-layer array.  Its maximum error is already bounded by
the preceding induction step.  A previous token's error is never
fed back into the same layer's residual state.  The only
cross-position operation is the normalized weighted average,
whose uniform Lipschitz bound was proved above.  Exact integer
moment sums evaluate the attention of the approximated arrays
without a further time-accumulation error.  The induction
therefore runs over the $J$ network stages only.  The explicit
polynomial bound
\eqref{eq:momentexplicitcost} completes the claim.
\end{proof}

\begin{lemma}[Exact sampling from certified replay]
\label{lem:replayexactsampling}
Suppose a deterministic causal evaluator can maintain a token
history at any fixed accuracy $2^{-p}$, spending at most
$G(p+\log(T+2))$ bits per appended token, where $G$ is a fixed
polynomial with nonnegative coefficients.  Suppose a replay of
all $T$ tokens costs at most $T G(p+\log(T+2))$ and returns
certified intervals of width at most $2^{-p}$ for all categorical
cumulative probabilities.  Then exact autoregressive sampling
can be maintained with expected polynomial work per token,
including precision changes, after polynomial prefill.
\end{lemma}

\paragraph{Proof idea.}
Routine precision is a sufficiently large multiple of $\log T$.
Only a uniform random number close to a cumulative boundary
requires a complete finer replay.  Background processing prepares
the next routine index before the current precision expires.
The interval decision is classical.  The additional point here
is paying for a complete causal replay only on its small
ambiguity event, while preparing routine precision increases
with bounded work at each token.

\begin{proof}
For a power-of-two horizon $R$, choose
\[
 p_R=\left\lceil\log_2\bigl(64V(R+2)^3\bigr)\right\rceil.
\]
At a query draw the first $p_R$ bits of a fresh uniform binary
expansion $U$.  Let $I_U=[a,b]$ be its dyadic interval, of width
$2^{-p_R}$.  Denote the certified cumulative intervals by
$[\underline F_i,\overline F_i]$, and use the exact endpoints
$F_0=0$ and $F_V=1$.  If
$\overline F_{i-1}\le a$ and $b\le\underline F_i$, output
token $i$.  Equality at a dyadic cell endpoint is permitted; the
rule can then differ from the inverse CDF only when $U$ equals a
true boundary, an event of probability zero.  Otherwise
replay the current history at accuracy $p_R+1$, reveal one more
uniform bit, and try again.  Continue with $p_R+j$ as needed.
Intervals from successive replays need not be nested; every one
contains the same true cumulative probabilities.

If the precision-$p$ comparison is unresolved, $U$ is within
$2^{1-p}$ of one of the at most $V-1$ cumulative boundaries.
The conditional probability of this event, given the token
history, is at most
\begin{equation}
 4V2^{-p}.
 \label{eq:replayuncertainty}
\end{equation}
Each realized uniform number outside the finite set of true
boundaries eventually lies in a unique certified categorical
interval.  The algorithm therefore stops almost surely, and its
output is exactly the inverse-CDF sample.  It does not substitute
an approximate probability law.

\paragraph{Cost of temporary replays.}
The expected work of temporary replays is bounded by
\[
 8VT2^{-p_R}
 \sum_{j\ge0}2^{-j}G(p_R+j+\log(T+2)).
\]
As long as $T\le R$, its prefactor is at most a constant times
$(R+2)^{-2}$.  If $G$ has degree $k$, the sum is polynomial
in $p_R+\log(T+2)$.  One explicit inequality, valid for $P\ge1$,
is
\[
 \sum_{j\ge0}2^{-j}(P+j)^k
 \le2^{k+1}k!P^k.
\]
It follows from $(P+j)^k\le P^k(j+1)^k$ and
$(j+1)^k\le k!\binom{j+k}{k}$, followed by the generating
function for these binomial coefficients.  This proves the
expected work bound without assuming finite expected stopping
time in advance.  The same tail bounds all additional random
bits.  After the token is returned, discard every temporary
fine-precision index and retain the deterministic routine index.

\paragraph{The background schedule.}
It remains to maintain routine precision without a large online
rebuild.  Let $N$ be the least power of two at least $T_0$.
During prefill build the routine index at precision $p_{2N}$.
This horizon is at least $2T_0$, so the first append is not forced
to pay for a prefix rebuild.  Before context length $N$, use
this index alone.  At length $N$, begin an empty future index at
precision $p_{4N}$.  During each of the next $N$ appends, perform
the current index's one-token update and replay two successive
history tokens into the future index.  After $j$ such appends,
the history has length $N+j$ and the future index has replayed
$2j$ tokens.  Since $2j\le N+j$ for $j\le N$, the next replay
tokens always exist.  At length $2N$ the future index has replayed
exactly the entire current history.  Switch to it, replace $N$
by $2N$, and begin the next future index.

Every online append therefore performs at most three deterministic
token evaluations.  The future precision differs from the
current precision by only a constant number of bits in $p_R$;
its other model-dependent guards are unchanged.  No operation
scans the context as part of the deterministic maintenance
schedule.  The only complete scans are temporary random
refinements, whose expected cost was bounded above.

\paragraph{The exact law.}
Store the actual token history for all replays.  A replay is
driven only by this history and the fixed model.  It never uses
the approximate hidden states of a different precision index.
Conditional on every realized token history, its intervals
therefore enclose the same exact next-token law.  Fresh uniform
bits give the correct next conditional distribution, and induction
establishes the entire autoregressive joint law.  Background
replays and temporary refinements do not change any previously
sampled token.
\end{proof}

\paragraph{Proof idea for Theorem~\ref{thm:momentdecoder}.}
The moment index supplies a maintained evaluator at each fixed
precision.  Uniform stability sets the required precision from
network depth, and the replay lemma handles its increase with
context length.  Combining their bit bounds charges both normal
updates and the rare refinements needed for the exact law.

\begin{proof}[Proof of Theorem~\ref{thm:momentdecoder}]
Apply Lemma~\ref{lem:prefixstability} to obtain the deterministic
fixed-precision evaluator, with the exact moment counters of
Lemma~\ref{lem:finitemomentindex}.  Its one-token work is bounded
by \eqref{eq:momentexplicitcost}, with
$m=O_r(\beta K^2+P+1)$ and the working precision stated in the
stability lemma.  Apply Lemma~\ref{lem:replayexactsampling}.
The explicit geometric-moment inequality converts the temporary
replay cost into another polynomial in the same parameters.
All storage consists of the fixed model, the token history, at
most two routine moment indices, and temporary replay storage
when required.  The routine indices have polynomial size in the
displayed parameters and $\log T$.  Temporary precision has no
deterministic upper bound, but the same geometric moment estimate
bounds its expected space by such a polynomial.  No infinite
expansion or real arithmetic oracle is used.

For the normalized model, $K=O(s\sqrt n)$, so
$H=O_r(1+\beta s^2n)$.  At $\beta=\sqrt n$ this is polynomial
in $n$ and $s$.  The normalization guard depends on $b$ only
through $\log L_0$, and the degree of the polynomial approximation
has no factor $2^b$.  Thus a very small
positive normalization floor changes precision polynomially.
The stipulated bounds on the remaining parameters give the
claimed $\operatorname{poly}_r(n,\log T)$ work.  Substituting
other polynomially evaluable updates only changes the range and
Lipschitz constants in the same finite-depth induction.
\end{proof}

\begin{corollary}[The affine output head needs no fixed logit range]
\label{cor:affineheadmoment}
The conclusions of Theorem~\ref{thm:momentdecoder} remain valid
when the promise that logits lie in $[-1,1]$ is removed, keeping
the stated affine-head row and bias bounds.  The cost is bounded
by another constructive polynomial in the same parameters.
The fixed-model corollaries below also hold for this extension.
\end{corollary}

\paragraph{Proof idea.}
Pre-normalization bounds every attention update, so additive
residuals give a polynomial logit range.  Subtracting an upper
bound for the largest logit makes all exponentials nonpositive
and leaves a constant lower bound on their total.  Certified
softmax evaluation therefore has no cost exponential in the
logit range.

\begin{proof}
With the power-of-two value bound $K$ used above, the residual
range after $D$ blocks is at most
\[
 R_D=s+D(sK+s+1).
\]
Indeed the embedding range is at most $s$; an attention mean has
range at most $K$, its optional affine output map has range at
most $sK+s$, and the tanh update has range at most one.  This
bound also covers convex residual coefficients.  Hence each
output logit has absolute value at most
$Z=sR_D+s$, which is polynomial in $n,D,s$.

For completeness, consider certified intervals of width at most
$1/4$ for the real logits $z_1,\ldots,z_V$.  Let $A$ be the
largest of their rational upper endpoints.  Then every
$z_i-A$ is nonpositive and at least one is at least $-1/4$.
Consequently
\[
 S=\sum_i e^{z_i-A}\ge e^{-1/4}>1/2.
\]
Enclose each exponential to width
$\eta=2^{-p}/(64V)$.  On the negative half-line the exponential
has Lipschitz constant at most one, so this only requires
$p+O(\log(V+2))$ certified logit bits, followed by the
negative-exponential routine.  Intersect any refined logit
interval with its original enclosure, so its upper endpoint
remains at most $A$.  It does not require an exact
comparison of two unknown logits.

Every numerator partial sum and the denominator has error at
most $V\eta$.  Put $\delta=V\eta\le S/2$.  Outward interval
division, after clipping each numerator interval to nonnegative
values, has width at most
\[
 \frac{N+\delta}{S-\delta}
 -\frac{N-\delta}{S+\delta}
 \le\frac{4S\delta}{S^2-\delta^2}
 \le\frac{16\delta}{3S}
 <11\delta<2^{-p-2}.
\]
Here $0\le N\le S$ is the true partial sum.  A final outward
dyadic rounding leaves
width below $2^{-p}$.  All exponents have polynomial encoding
length and are nonpositive; no factor $e^Z$ occurs in their
bit cost.  The global score-to-probability Lipschitz bound for
softmax used in Lemma~\ref{lem:prefixstability} is valid at
every logit range.  Thus both its stability guard and the replay
argument apply with the stated polynomial evaluation cost.
\end{proof}

\begin{corollary}[Specified GELU and SwiGLU updates]
\label{cor:momentgeluswiglu}
Theorem~\ref{thm:momentdecoder}, with the unrestricted affine
head of Corollary~\ref{cor:affineheadmoment}, also holds when
the tanh feedforward update is replaced by either exact GELU
or SwiGLU applied to affine projections of the normalized
stream, followed by an affine output map.  The hidden width
is included as a polynomial parameter, or is bounded by a
fixed polynomial in $n$.  All affine row and bias bounds
remain at most $s$.  The update can be multiplied by any fixed
dyadic coefficient $\alpha\in[0,1]$.
\end{corollary}

\paragraph{Proof idea.}
Normalization bounds each affine input independently of the
current residual magnitude.  GELU and SwiGLU consequently
have polynomial range and derivative bounds.  Their scalar
functions have certified finite series evaluations, so they
satisfy the explicit hypotheses of the decoder theorem.

\begin{proof}
Write $K=s(\sqrt n+1)$, and increase it to an integer when
forming series bounds.  Every coordinate of every affine
projection of a normalized vector has magnitude at most $K$.
For exact GELU, put $g(u)=u\Phi(u)$, where $\Phi$ is the
standard normal distribution function and $\phi$ its density.
Since $0\le\Phi\le1$ and $0\le\phi\le1$,
\[
 |g(u)|\le K,\qquad
 |g'(u)|=|\Phi(u)+u\phi(u)|\le1+K
 \quad (|u|\le K).
\]
For a SwiGLU coordinate put $w(u,v)=u v\sigma(v)$, where
$\sigma(v)=(1+e^{-v})^{-1}$.  The identities
$\sigma'=\sigma(1-\sigma)$ and $0\le\sigma'\le1/4$ give
\[
 |w(u,v)|\le K^2,\quad
 |\partial_u w|\le K,\quad
 |\partial_v w|\le K(1+K/4).
\]
Thus its Lipschitz constant from the two-coordinate maximum
norm is at most $2K+K^2/4$.  Composing the input projections
and the output map multiplies these bounds by at most $s^2$.
The output ranges are at most $sK+s$ and $sK^2+s$,
respectively.  These estimates do not depend on hidden width
because they use row sums.  Multiplication by
$\alpha\in[0,1]$ only decreases both range and derivative bounds.
Adding $D$ such bounded updates leaves a residual and affine
head range polynomial in $n,D,s$, as required by
Corollary~\ref{cor:affineheadmoment}.

Here is a direct certified evaluation procedure.  The sigmoid
uses a negative exponential and a denominator at least one:
for $v\ge0$ use $1/(1+e^{-v})$, and for $v<0$ use
$e^v/(1+e^v)$.  Evaluation at a dyadic argument therefore has
polynomial cost in $K$ and the requested precision, by
Lemma~\ref{lem:newnegativeexp}.  For GELU use
\[
 \Phi(u)=\frac12+\frac1{\sqrt{2\pi}}
                  \int_0^u e^{-t^2/2}\,dt.
\]
Expand the integrand in powers of $-t^2/2$ and integrate its
finite polynomial exactly.  Put $H=\max(1,K^2/2)$.  The same
factorial remainder bound as in
\eqref{eq:momenttaylorerror} shows that
$O(H+p+\log(K+2))$ terms suffice for integral error
$2^{-p-O(1)}$ uniformly on $[-K,K]$.
The intermediate rational coefficients and arithmetic have
polynomial bit length in $K,p,b$.  Certified values of
$\pi$ follow from the arctangent identity used below in
Corollary~\ref{cor:fixedpositionalrotation}, and rational
bisection gives its square root and reciprocal.  Guard
precision $O(K^2+p+\log(K+2))$ controls every product and
final division.  Request an additional $\lceil\log_2(K+2)\rceil$ bits for
$\Phi$ before multiplying it by $u$.  Thus both updates have
polynomial certified evaluation cost.  For a real input already enclosed by a
certified interval, evaluate at its dyadic midpoint and add
the displayed derivative bound times its radius.  No exact
comparison with a transcendental value is required.

The range bounds, evaluation algorithms, and derivative
bounds verify the specified-family assumptions of
Theorem~\ref{thm:momentdecoder}.  Its depth stability and
replay proofs therefore apply without further changes.
\end{proof}

\begin{corollary}[A completely fixed model]
\label{cor:fixedmodelcontext}
Fix every parameter of one finite model covered by
Theorem~\ref{thm:momentdecoder}, including its width, depth,
vocabulary and weights.  Its key ranks are then fixed integers,
even if its key maps have full rank.  After prefill, exact
autoregressive decoding takes $\log^{O_{\mathrm{model}}(1)}(T+2)$ expected bit
work per token.  For the basic tanh model, including the affine
head extension, the explicit implementation above permits the
bound $O_{\mathrm{model}}(\log^{2r+12}(T+2))$.  Other specified
feedforward families retain the $\log^{O_{\mathrm{model}}(1)}(T+2)$
bound, whose exponent includes their evaluation polynomial.
\end{corollary}

\paragraph{Proof idea.}
Only the requested precision increases with context length.
The number of monomials is a polynomial in that precision when
the key dimension is fixed.

\begin{proof}
For a fixed model, $H$ and all stability constants are constant.
Routine accuracy and working precision are $O_{\mathrm{model}}
(\log(T+2))$, as is $m$.  The integer length $L$ in
\eqref{eq:momentexplicitcost} is
$O_{\mathrm{model}}(\log^2(T+2))$.  For the basic tanh model, substitution gives exponent
$2r+12$, and the temporary-replay geometric moment bound
preserves that degree.  The background schedule has constant
overhead.  Taking the largest key rank across layers proves the
claim.  For another specified feedforward family, replace this
explicit polynomial by the fixed polynomial that also charges
its certified scalar evaluations; the same argument preserves
a model-dependent polylogarithmic exponent.
\end{proof}

\begin{corollary}[Fixed positional rotations]
\label{cor:fixedpositionalrotation}
In the completely fixed model of
Corollary~\ref{cor:fixedmodelcontext}, allow query and key vectors
at position $t$ to undergo coordinate-pair rotations with angles
$t\omega_1,\ldots,t\omega_{\lfloor n/2\rfloor}$.
Each frequency is fixed and comes with a specified finite
algorithm that returns a certified interval of width $2^{-p}$
in time polynomial in $p$.  Fixed rational and real algebraic
frequencies satisfy this assumption.  Exact autoregressive
decoding still has $\log^{O_{\mathrm{model}}(1)}(T+2)$ expected
bit work at each token after prefill.
The same conclusion holds for a fixed finite collection of such
rotations and fixed unrotated coordinates.
\end{corollary}

\paragraph{Proof idea.}
Rotation does not increase Euclidean length and has a bounded
coordinate Lipschitz factor.  Evaluating an angle at a late
position needs additional frequency bits proportional to
$\log T$.  The full ambient key space avoids a rank condition
on the rotating subspaces.

\begin{proof}
A pair rotation has matrix
\[
 \begin{pmatrix}\cos\theta&-\sin\theta\\
                 \sin\theta& \cos\theta\end{pmatrix}.
\]
Its maximum absolute row sum is at most $\sqrt2$, uniformly in
$\theta$.  Consequently it introduces only a fixed factor in
the key and query coordinate bounds and in their Lipschitz
bounds with respect to preceding-layer states.  The position
is a known exact integer and is never perturbed.

We give a finite computation of its coefficients.  Fix a
frequency bound $C_\omega\ge1+|\omega|$.  At position
$1\le t\le T$, form a dyadic approximation $\theta_0$ to
$t\omega$, with error at most $2^{-p-10}$, by requesting
$p+\lceil\log_2(T+1)\rceil+O(1)$ frequency bits.
Compute a dyadic $\pi_0\in[3,4]$ with error at most
$2^{-p-10}/[C_\omega(T+1)]$.
This has polynomial bit cost in $p+\log(T+2)$: for example,
the identity
\[
 \pi=16\arctan(1/5)-4\arctan(1/239)
\]
and the alternating arctangent series give certified intervals
using a linear number of terms in the requested precision,
with polynomial-length integer arithmetic.
Set $k=\lfloor\theta_0/(2\pi_0)\rfloor$ by exact rational
division and put $u_0=\theta_0-2k\pi_0$.
Then $0\le u_0<8$ and $|k|\le C_\omega(T+1)$.
The true reduced argument $u=t\omega-2k\pi$ satisfies
\[
 |u-u_0|\le|t\omega-\theta_0|+2|k||\pi-\pi_0|
          \le3\cdot2^{-p-10}.
\]
No test of whether a true angle lies exactly on a period
boundary is made.  Sine and cosine have Lipschitz constant one.
Their rational Taylor polynomials through degree $32(p+8)$
on $[0,8]$ have remainder at most
$8^{32(p+8)+1}/[32(p+8)+1]!<2^{-p-6}$.
Evaluating these finite polynomials exactly and rounding outward
gives certified error $2^{-p-4}$ in polynomial bit work.
Thus both coefficients have the
required certified evaluation bound.  For a fixed algebraic
frequency, an isolating interval for its minimal polynomial
and rational bisection give the assumed polynomial-time
frequency approximations.

Use the common ambient chart $A=I_n$, $c_0=0$, of dimension
$r=n$.  After a rotation, round each approximate query and key
coordinate to the common working grid by the certified midpoint
rule.  The keys and the coefficient vector $t$ then have the finite
form that Lemma~\ref{lem:finitemomentindex} requires.  At full
dimension, rounding the keys cannot violate the key-space condition.
Request an additional constant number of coefficient bits to
account for the fixed coordinate range.  The rotation error
and the final grid rounding are then included in the same
local error budget $\delta$.  The resulting moment counters
remain exact integers.  All score and state stability bounds
are independent of position, while coefficient evaluation
adds only a fixed polynomial in $p+\log(T+2)$ to per-token
work.  Lemma~\ref{lem:replayexactsampling} applies with this
enlarged polynomial and proves the claimed conditional exact
law and expected cost.
\end{proof}

\paragraph{Scope.}
For a completely fixed full-rank model,
Corollaries~\ref{cor:fixedmodelcontext}
and~\ref{cor:fixedpositionalrotation} are formal asymptotic
statements with a potentially enormous model-dependent exponent and
constant.  They do not assert a practical improvement at present
model dimensions.  The moment dimension grows as
$\binom{O_r(\beta s^2n+P)+r}{r}$.
The source of the positive result is a structural restriction on
the key maps.  The arbitrary-fresh-query attention hardness results
of Section~\ref{sec:indexedlower} permit ranks that grow with the
context, so they do not exclude this theorem.  Conversely, the theorem does not remove their
barrier for unrestricted caches.

